\chapter{Implementation and Module-Wise Coding}\label{ch:impl}

\section{Development Environment}
The environment is given in Table~\ref{tab:env}. Everything runs on a single Windows 11
laptop. The backend and the portal are ordinary Maven projects that produce executable
JAR files; MySQL runs as a Windows service. Maven and the JDK are the only build tools; no
Node.js toolchain is part of the product (Node.js was used only to capture screenshots
for this report).

The three source folders of the repository are:
\begin{itemize}
\item \cd{database/} --- the SQL scripts, seed data, test script and advanced queries;
\item \cd{backend-java/} --- the Spring Boot REST backend (152 main source files, 15 integration-test classes);
\item \cd{frontend-java/} --- the Spring Boot portal (18 controllers, 70 templates).
\end{itemize}

\section{Overall Implementation Strategy}
The work followed the dependency order of the three tiers. The database came first
because the backend and the portal both depend on its table and column names. The backend
came next because it defines the API that the portal calls. The portal was built on top, in
two stages. The database was then hardened with constraints, indexes, views, triggers and
procedures, after the backend existed, so that the triggers could be designed around what
the Java services already write. The last stage was validation: running the whole test set,
crawling the portal, measuring performance, and fixing what was found
(Chapter~\ref{ch:test}).

\section{Module-Wise Implementation}
Every module has the same shape: a controller that declares the URL and the permission,
DTO records that validate the input, a service that holds the rules inside a transaction,
and entities that map the tables. The paragraphs below give the implementation points that
are specific to each module.

\paragraph{Authentication and access control.}
\cd{AuthService} implements sign-in, refresh and sign-out; \cd{JwtService} creates and
verifies HS256 tokens; \cd{UserLoader} builds the \cd{CurrentUser} of each request from
the database. \cd{RbacService} lets an administrator create roles and users, set a role's
permissions and assign roles. Passwords are hashed with BCrypt at cost 12. At first start
a command-line step creates the roles, the 62 permissions and the first administrator.

\paragraph{Citizens, staff and offices.}
A citizen is created together with a person row and an optional first address in one
transaction. A new address closes the previous current address of the same type instead of
overwriting it, so the address history is kept. Employees have dated postings, and a
generated column permits only one open posting per employee.

\paragraph{Applications.}
\cd{ApplicationService} generates the application number, computes the fee as the service's
base fee plus every fee component in force on that date, and moves the status through the
transition table. Approval is refused unless the required documents are verified and, when a
fee is due, a successful payment exists. Every move writes a history row and an audit row.

\paragraph{Documents and appointments.}
\cd{DocumentService} checks the extension, the leading bytes and the size of an upload,
stores the file under the storage directory and keeps only its path and metadata in the
database. A verified document can expire after the number of days defined by its document
type. \cd{AppointmentService} creates slots (refusing duplicates, which the unique key
cannot catch when the counter is empty because MySQL treats NULLs as distinct), and books,
cancels and reschedules seats under the locking rules of Chapter~\ref{ch:algo}.

\paragraph{Licensing.}
\cd{LicenceService} issues learner licences (one active per citizen, enforced by a generated
column), records driving tests, issues driving licences with one or more classes, and
suspends, revokes, reinstates and renews them within the transition table. An expiry
sweep closes licences whose date has passed. A licence can require a passed test when tests
exist for the citizen.

\paragraph{Vehicles, ownership and compliance.}
\cd{VehicleService} normalises registration numbers, rejects duplicate registration,
chassis and engine numbers and a model that does not belong to the manufacturer, and
creates the first ownership row. \cd{OwnershipService} implements the transfer algorithm.
\cd{ComplianceService} records inspections, fitness, pollution, insurance and road-tax
records on active vehicles, answers the ``what expires soon'' question with a combined
query, and refreshes time-driven statuses.

\paragraph{Permits, violations and challans.}
\cd{PermitService} manages routes with ordered segments and permits whose expiry defaults from
the permit type. \cd{ViolationService} records violations against vehicles (the driver may be
unknown), bundles them into a challan whose total is the exact decimal sum, and drives the
challan state machine, including the return to \texttt{ISSUED} when a dispute is dismissed.

\paragraph{Payments.}
\cd{PaymentService} implements the idempotent payment algorithm, partial and full refunds,
and the rule that only one refund may be in flight at a time. A refund of a road-tax
payment returns the road-tax record to its previous state.

\paragraph{Support, notifications and audit.}
\cd{SupportService} handles complaints and appeals; an upheld appeal acts on its target
(a challan is cancelled, a rejected application is reopened, a suspended licence is
reinstated). \cd{Notifications} queues messages inside the business transaction and delivers
them after it commits. \cd{Audit} writes the audit row. \cd{DashboardService} returns the
live counts shown on the dashboard.

\section{Front-End Implementation}
\subsection{Structure}
Each portal controller does three things: it reads request parameters, calls the backend
through \cd{ApiClient}, and passes the resulting JSON to a Thymeleaf template. Thymeleaf
reads the JSON directly with expressions such as \cd{row.path('citizen_code').asText()}, so the
portal has no duplicate model classes to keep in step with the backend.

\subsection{The API client}
\cd{ApiClient} adds the session's bearer token to each call. If the backend answers 401 it
tries once to refresh the token pair; if that fails the session is cleared and the user is
sent to the sign-in page. Error bodies are converted to one exception type carrying the
status, the code, the message and the list of field errors. Before a request is sent, the
paging parameter is clamped to the range 1 to $2^{31}-1$ and a search term longer than 100
characters is cut to 100, because the API would otherwise reject a hand-edited address
with a 422; this keeps such addresses working.

\subsection{Templates}
All pages share a layout fragment with the sidebar and a flash-message area, and a small
library of fragments for text inputs, selects, pagers and status badges. The sidebar shows an
entry only if the session holds the matching permission: an administrator sees 27 entries,
a citizen sees five, including the dashboard (Figure~\ref{fig:citizenmenu}). Forms post to the portal, which sends a
compact JSON body to the API (blank optional fields are omitted when creating and sent as
\texttt{null} when updating so that a field can be cleared); a validation error is mapped back
to the field that caused it and the form is shown again with the user's entries.

\subsection{Screens}
Figures~\ref{fig:scrdash}, \ref{fig:scrapp}, \ref{fig:scrveh}, \ref{fig:scrrole},
\ref{fig:citizenmenu} and~\ref{fig:scrval} show six representative screens captured from
the running system with the demonstration data. They are discussed with the results in
Chapter~\ref{ch:results}.

\section{Back-End Implementation}
\subsection{Request pipeline and conventions}
The pipeline is shown in Figure~\ref{fig:request}. A few conventions run through the API:
\begin{itemize}
\item JSON field names are \texttt{snake\_case}; dates are ISO-8601; money is a string with two decimals (\texttt{"1500.50"}), never a floating-point number;
\item a list returns \texttt{items}, \texttt{page}, \texttt{page\_size}, \texttt{total} and \texttt{pages}, with at most 100 items per page and a whitelist of sort keys;
\item the database session uses UTC, and times are converted for display only;
\item the OpenAPI description is generated by springdoc and served as Swagger UI at \cd{/docs}.
\end{itemize}

\subsection{Persistence}
Entities are plain mappings of tables with foreign keys held as ids rather than as object
references, which avoids lazy-loading surprises and makes the locking calls explicit. The
generated columns are mapped read-only. A tiny helper, \cd{Db}, wraps queries, row locks
(each lock refreshes the entity so that nothing read before the lock is stale), paging and
native queries. A schema-audit test verifies that every table and column has a mapping.

\section{Database Implementation}
\subsection{Scripts and order}
The scripts are numbered in the order they are run: files 01 to 14 create the database and
the 63 tables; 15 adds CHECK constraints; 16 adds indexes; 17 creates views; 18 creates
triggers; 20 creates procedures; 21 holds the analytical queries; 22 and 23 are seed data;
the verification folder holds the SQL test script. The file \cd{00_run_all.sql} runs the
creation scripts with \texttt{SOURCE} in this order.

\subsection{Seed data}
The first seed (file 22) is small, with 10 citizens, 10 applications and 5 vehicles. A
second, comprehensive seed is generated by a Python script and holds 23 citizens, 20
applications, 15 vehicles, 11 driving licences, 13 payments and 61 audit entries. It covers
\emph{every value of every ENUM column}: we checked this by comparing the allowed values of
all 25 ENUM columns with the values present. It includes edge cases such as a person with
accents and an apostrophe in the name, a person of gender \texttt{OTHER}, citizens with no
activity at all, and a long audit trail for pagination. All event dates are written as
offsets from the current date, so ``expires in four days'' remains true whenever the seed
is loaded. A foreign-key audit of the loaded data found and removed four orphan references,
which the seed's \texttt{FOREIGN\_KEY\_CHECKS = 0} setting had hidden (Section~\ref{sec:defects}).

\subsection{Environment and accounts}\label{sec:envsetup}
The application does not connect as the MySQL administrator. A dedicated account has rights
only on the project's two databases. Table~\ref{tab:setup} lists the set-up steps that were
used on the development machine; passwords and the signing secret are placed in the
git-ignored \texttt{.env} file and are not reproduced here.

\begin{table}[H]
\centering\small
\caption{Set-up steps (secrets omitted)}\label{tab:setup}
\begin{tabularx}{\linewidth}{@{}p{0.8cm}X@{}}
\toprule \textbf{Step} & \textbf{Action}\\ \midrule
1 & Install JDK 21 and Maven 3.9; confirm with \cmd{java -version} and \cmd{mvn -v}\\
2 & Create the schema: run \cmd{00_run_all.sql} with the mysql client from the \texttt{database} folder\\
3 & Create the application account with \texttt{GRANT ALL} on the project database and on the test database only\\
4 & Create \texttt{backend-java/.env} from the example: database host, port, user, password; a random JWT secret of at least 32 characters; the first administrator's username and password\\
5 & Build: \cmd{mvn -DskipTests package} in each of the two projects\\
6 & Bootstrap: \cmd{java -jar target/rto-management-backend-1.0.0.jar --rto.command=bootstrap}\\
7 & Start the backend, then the portal; open \cmd{http://localhost:8080}\\
\bottomrule
\end{tabularx}
\end{table}

\section{Communication and Concurrency Implementation}
The portal and backend communicate over HTTP with JSON bodies. The portal is
stateful (a session per browser), the backend is stateless (everything it needs is in the
token and the database). The backend's connection pool has at most 20 connections, and the
connections use the \texttt{READ COMMITTED} isolation level. Concurrency control is done with
row locks and guarded updates inside each transaction, as described in
Chapter~\ref{ch:algo}. The tests that exercise it start several threads together using a
latch and then check the resulting data (Section~\ref{sec:conc}).

\section{Git Repository and Timeline}\label{sec:timeline}
The source is kept in a Git repository (\texttt{github.com/simritt/rto-management}).
Table~\ref{tab:timeline} gives the order of the work from the commit history. Dates are
commit dates, which show when work was pushed, not necessarily when it was written.

\begin{table}[H]
\centering\small
\caption{Development timeline from the repository history}\label{tab:timeline}
\begin{tabularx}{\linewidth}{@{}p{2.3cm}p{4.4cm}X@{}}
\toprule \textbf{Date} & \textbf{Work} & \textbf{Reason for the order}\\ \midrule
10 Aug 2026 & Database schema: 63 tables, run-all and verification scripts & Both other tiers depend on its names and keys\\
4 Oct 2026 & Spring Boot backend with RBAC, transactions, locking, audit, notifications and integration tests & Needs only the schema; defines the API\\
7 Oct 2026 & First version of the portal (feature branch) & Needs a working API; kept off the main branch until ready\\
8 Oct 2026 & Eleven commits: constraints, indexes, views, triggers, locking examples, procedures, advanced queries, seed, run-all update, SQL tests, \texttt{.gitignore} & Added after the backend so triggers avoid duplicating its writes\\
9 Oct 2026 & Project README & Written once the structure and run steps were settled\\
20 Oct 2026 & Complete portal modules and workflows; merge into the main branch & Last, because it depends on every other layer\\
22 Oct 2026 & Fix: challan insert violating the positive-amount rule; real BCrypt hashes for seed users & Corrections to the seed data\\
3 Nov 2026 & Portal error-handling and input hardening; comprehensive demonstration seed & Result of the final validation (Chapter~\ref{ch:test})\\
\bottomrule
\end{tabularx}
\end{table}
The database enhancements were committed one concern per commit, in the order in which a
database is hardened: constraints to define what is valid, indexes to make valid data
quick to find, views for simple reading, triggers for automatic reactions, transactions
and locking for correctness under concurrency, procedures to package multi-step
operations, queries to show analysis, and seed data and tests to demonstrate and verify. A
reviewer can therefore read or revert one concern without touching the rest. The commit
messages follow the conventional style (\texttt{feat(db):}, \texttt{test(db):},
\texttt{docs:}, \texttt{fix(frontend):}).

\section{Integration of Modules}
A typical end-to-end scenario touches most modules. A clerk registers a citizen and creates
an application; the citizen uploads a document; a clerk verifies it; a slot is booked;
the fee is paid; an officer approves; a licence is issued; the audit log shows each step;
and notifications are delivered after each commit. The integration tests run such
flows (for example \texttt{fullHappyPathRecordsEveryStep}) against a real database, and the
portal crawl exercises the same pages a user would visit.

\section{Implementation Challenges and Resolutions}
Table~\ref{tab:challenges} summarises the implementation-level problems we met and how
they were resolved. The defects found during final validation are described separately in
Section~\ref{sec:defects}.

\begin{table}[H]
\centering\small
\caption{Implementation challenges and resolutions}\label{tab:challenges}
\begin{tabularx}{\linewidth}{@{}p{4.3cm}X@{}}
\toprule \textbf{Challenge} & \textbf{Resolution}\\ \midrule
Triggers and Java both writing history rows & Analysed what each service writes; added triggers only for licence and permit history and for delete guards\\
New ownership row rejected by the unique key & Close the old row and flush it to the database before inserting the new row\\
Duplicate slots when the counter is empty & MySQL treats NULLs as distinct in a unique key, so the service checks for duplicates under an office row lock\\
Possible deadlock between two bookings & Always lock slots in ascending id order; fixed lock order application, appointment, slot\\
Money shown with rounding errors & \texttt{DECIMAL} columns, \texttt{BigDecimal} in Java, strings in JSON\\
Time-zone drift between server, JVM and database & The database session and the JVM both use UTC; the connection URL sets the connection time zone\\
Token with junk after the signature accepted by a lenient decoder & Require a canonical three-part token and compare re-encoded parts\\
A repeated payment request returned an error instead of the first result & Answer the replay before validating the target\\
Triggers and procedures could not be created by the application account & With binary logging on, MySQL needs administrator rights or the \texttt{log\_bin\_trust\_function\_creators} setting before a non-administrator can create triggers or stored routines; the administrator enabled the setting\\
Thymeleaf variable named \texttt{session} clashed with its built-in & Renamed the variable to \texttt{user}\\
\bottomrule
\end{tabularx}
\end{table}
